% Part: lambda-calculus
% Chapter: syntax
% Section: unique-readability

\documentclass[../../../include/open-logic-section]{subfiles}

\begin{document}

\olfileid{lam}{syn}{unq}
\olsection{خوانش یکتا}

ممکن است بپرسیم آیا هر ترم فقط به یک شیوه ساخته می‌شود؛ پاسخ مثبت است.
هر ترم لامبدا تنها یک شیوهٔ ساخت و تفسیر دارد. در بحث زیر، \emph{ساخت}
یعنی رویهٔ ساختن یک ترم با استفادهٔ یک‌باره یا چندباره از قواعد ساخت در
\olref[trm]{defn:term}.

\begin{lem}\ollabel{lem:term-start}
هر ترم یا با یک متغیر آغاز می‌شود یا با یک پرانتز.
\end{lem}

\begin{proof}
چیزی تنها هنگامی ترم به شمار می‌آید که مطابق~\olref[trm]{defn:term}
ساخته شود. اگر حاصل~\olref[trm]{defn:term-var} باشد، باید متغیر باشد. اگر
حاصل~\olref[trm]{defn:term-abs} یا~\olref[trm]{defn:term-app} باشد، با
پرانتز آغاز می‌شود.
\end{proof}

\begin{lem}\ollabel{lem:app-start}
حاصل یک کاربرد یا با دو پرانتز آغاز می‌شود یا با یک پرانتز و یک متغیر.
\end{lem}

\begin{proof}
اگر~$M$ حاصل یک کاربرد باشد، به‌شکل~$(PQ)$ است و بنابراین با پرانتز آغاز
می‌شود. چون~$P$ یک ترم است، بنا بر~\olref{lem:term-start} یا با پرانتز
آغاز می‌شود یا با متغیر.
\end{proof}

\begin{lem}\ollabel{lem:initial}
هیچ بخش آغازین سره‌ای از یک ترم، خود یک ترم نیست.
\end{lem}

\begin{prob}
با استقرا بر طول ترم‌ها، \olref[lam][syn][unq]{lem:initial} را ثابت کنید.
\end{prob}

\begin{prop}[خوانش یکتا] \ollabel{prop:unq}
برای هر ترم، ساختی یکتا وجود دارد. به بیان دیگر، اگر ترم~$M$ با ساختی
تشکیل شود، همان ساخت یگانه ساختی است که می‌تواند این ترم را تشکیل دهد.
\end{prop}

\begin{proof}
  این حکم را با استقرا بر ساخت ترم‌ها ثابت می‌کنیم.
  \begin{enumerate}
    \item $M$ به‌شکل~$x$ است، که~$x$ متغیری است. چون حاصل‌های انتزاع و
      کاربرد همواره با پرانتز آغاز می‌شوند، نمی‌توانسته‌اند در ساخت~$M$
      به کار رفته باشند؛ پس ساخت~$M$ باید تنها یک گام از
      \olref[trm]{defn:term}\olref[trm]{defn:term-var} باشد.
    \item $M$ به‌شکل~$(\lambd[x][N])$ است، که~$x$ متغیر و~$N$ ترم است.
      این ترم نمی‌تواند مطابق
      \olref[trm]{defn:term}\olref[trm]{defn:term-var} ساخته شده باشد،
      زیرا یک متغیر تنها نیست. بنا بر~\olref{lem:app-start} نیز حاصل یک
      کاربرد نیست. پس~$M$ فقط می‌تواند حاصل انتزاع بر~$N$ باشد. بنا بر
      فرض استقرا می‌دانیم که ساخت~$N$ نیز یکتاست.
    \item $M$ به‌شکل~$(PQ)$ است، که~$P$ و~$Q$ ترم‌اند. چون با پرانتز
      آغاز می‌شود، نمی‌تواند هم‌زمان با
      \olref[trm]{defn:term}\olref[trm]{defn:term-var} ساخته شده باشد.
      بنا بر~\olref{lem:term-start}، $P$ نمی‌تواند با~$\lambd$ آغاز شود؛
      پس~$(PQ)$ نمی‌تواند حاصل یک انتزاع باشد. اکنون فرض کنید شیوهٔ دیگری
      برای ساخت~$M$ با کاربرد وجود داشته باشد؛ مثلاً~$M$ همچنین به‌شکل
      $(P'Q')$ باشد. اگر~$P \neq P'$، یکی از~$P$ و~$P'$ بخش آغازین سرهٔ
      دیگری است، و این بنا بر~\olref{lem:initial} ناممکن است. پس~$P=P'$
      و در نتیجه~$Q=Q'$. بنابراین~$P$ و~$Q$ به‌طور یکتا تعیین می‌شوند و
      بنا بر فرض استقرا، ساخت‌های~$P$ و~$Q$ یکتا هستند.
  \end{enumerate}
\end{proof}

بازگویی خواناتر قضیهٔ بالا چنین است:
\begin{prop}
  ترم~$M$ تنها می‌تواند یکی از شکل‌های زیر را داشته باشد:
  \begin{enumerate}
    \item $x$، که~$x$ متغیری است که~$M$ آن را به‌طور یکتا تعیین می‌کند.
    \item $(\lambd[x][N])$، که~$x$ متغیر و~$N$ ترمی دیگر است و~$M$ هر
      دوی آن‌ها را به‌طور یکتا تعیین می‌کند.
    \item $(PQ)$، که~$P$ و~$Q$ دو ترم‌اند که~$M$ آن‌ها را به‌طور یکتا
      تعیین می‌کند.
  \end{enumerate}
\end{prop}

\end{document}
